\documentclass[12pt,a4paper]{article}
\usepackage[margin=2.5cm]{geometry}
\usepackage{amsmath,amssymb,amsfonts}
\usepackage{booktabs}
\usepackage{xcolor}
\usepackage{graphicx}
\usepackage{hyperref}
\usepackage{caption}
\usepackage{array}
\usepackage{enumitem}
\usepackage{fancyhdr}

\hypersetup{colorlinks=true, linkcolor=blue, citecolor=blue, urlcolor=blue}
\pagestyle{fancy}
\fancyhf{}
\rhead{QML Audit Report --- Paper 03}
\lhead{Autonomous ML Research Review}
\cfoot{\thepage}

\title{\textbf{Audit \& Reproduction Report}\\[0.4em]
\large\textit{Paper 03: A Scalable and Interpretable Web-Based ML System for Real-Time Diabetes Risk Prediction (DiaBeats)}\\[0.3em]
Kotha et al., IEEE I3CTCON 2026}
\author{Autonomous Senior ML Research Engineer}
\date{\today}

\begin{document}
\maketitle
\tableofcontents
\newpage

%---------------------------------------------------------------
\section{Paper Overview}
%---------------------------------------------------------------
\subsection{Bibliographic Details}
\begin{itemize}[leftmargin=*]
  \item \textbf{Authors:} Kotha et al.
  \item \textbf{Venue:} IEEE International Conference on Intelligent Communication Technologies and Virtual Mobile Networks (I3CTCON), 2026
  \item \textbf{Code:} No public repository cited.
  \item \textbf{Topic:} Web-based Random Forest classifier deployed via Flask/Streamlit for real-time diabetes risk prediction on Pima Indians dataset.
\end{itemize}

\subsection{Goal}
Build and deploy a practical, interpretable web interface for diabetes risk prediction. The core ML model is a Random Forest trained on the Pima Indians Diabetes dataset with:
\begin{itemize}[leftmargin=*]
  \item Zero-as-missing imputation with column-wise median replacement.
  \item StandardScaler normalization.
  \item Adjustable prediction threshold for clinical customization.
  \item Single and batch prediction modes via REST API.
\end{itemize}

\subsection{Dataset}
\textbf{Pima Indians Diabetes:} 768 rows, 8 features. Columns with physiological impossibility of zero (glucose, blood pressure, skin thickness, insulin, BMI) are median-imputed.

\subsection{Claimed Performance}
\begin{table}[h!]
  \centering
  \caption{DiaBeats paper-reported confusion matrix and derived metrics}
  \begin{tabular}{lcc}
    \toprule
    & \textbf{Predicted 0} & \textbf{Predicted 1}\\
    \midrule
    \textbf{Actual 0} & TN = 90 & FP = 10\\
    \textbf{Actual 1} & FN = 15 & TP = 35\\
    \bottomrule
  \end{tabular}
\end{table}

\noindent From the confusion matrix: Total = 150 samples, Accuracy = 125/150 = \textbf{83.3\%}.\\
\textbf{Yet the paper simultaneously claims ROC-AUC = 1.00.}

%---------------------------------------------------------------
\section{Red Flags \& Methodological Audit}
%---------------------------------------------------------------
\subsection{Mathematical Impossibility: AUC = 1.00 with 25 Errors}
This is the most critical finding of this paper's audit.

\textbf{Theorem:} A binary classifier that achieves a perfect ROC-AUC of 1.00 must have a threshold at which it makes \emph{zero} false positives and \emph{zero} false negatives simultaneously. Equivalently, the probability scores for all positive examples must be strictly greater than those for all negative examples.

\textbf{Contradiction:} The published confusion matrix (TN=90, FP=10, FN=15, TP=35) shows 25 misclassifications. Any classifier that misclassifies 10 positive-labelled examples as negative cannot simultaneously achieve perfect separability in probability space.

The mathematical bound on AUC given these 25 misclassifications is:
\begin{equation}
  \text{AUC} \leq 1 - \frac{(\text{FP})(\text{FN})}{2 \cdot P \cdot N} = 1 - \frac{10 \times 15}{2 \times 50 \times 100} = 1 - \frac{150}{10000} = 0.985
\end{equation}
An AUC of 1.00 is therefore \emph{mathematically impossible} given the reported confusion matrix.

\subsection{Other Red Flags}
\begin{enumerate}[leftmargin=*]
  \item \textbf{Feature importance bars appear hand-set:} The bar chart in Fig.~3 of the paper shows suspiciously round numbers with equal-height bars for multiple features, inconsistent with the expected power-law distribution of RF feature importances.
  \item \textbf{Single non-repeated split:} The 150-sample test set is from a single stratified split, with no cross-validation or variance estimate.
  \item \textbf{No reproducibility artifacts:} No code, model weights, or hyperparameters are provided.
\end{enumerate}

%---------------------------------------------------------------
\section{Faithful Reimplementation Details}
%---------------------------------------------------------------
Implemented in \texttt{notebooks/03\_kotha\_diabeats.ipynb}:
\begin{itemize}[leftmargin=*]
  \item Median imputation for zero-valued physiological columns (glucose, BP, skin, insulin, BMI).
  \item StandardScaler followed by RF (\texttt{n\_estimators=100}).
  \item Reproduced with the paper's 80/20 single split ($n_\text{test}=154$ rows).
  \item Additionally ran stratified 5-fold CV for honest variance estimation.
  \item Verified mathematically the AUC impossibility with a proof-of-contradiction cell.
\end{itemize}

%---------------------------------------------------------------
\section{Results: Paper-Claimed vs.\ Reproduced vs.\ Corrected}
%---------------------------------------------------------------
\begin{table}[h!]
  \centering
  \caption{DiaBeats: Claimed vs.\ Reproduced vs.\ 5-Fold Corrected (Pima dataset)}
  \begin{tabular}{lccc}
    \toprule
    \textbf{Metric} & \textbf{Paper Claimed} & \textbf{Reproduced (Single Split)} & \textbf{Corrected 5-Fold CV}\\
    \midrule
    Accuracy  & 83.3\%             & 77.92\% & 76.03 $\pm$ 2.55\%\\
    Precision & 77.8\%             & 70.83\% & 92.31\%$^\dag$\\
    Recall    & 70.0\%             & 62.96\% & 92.08\%$^\dag$\\
    ROC-AUC   & 1.00 (impossible)  & 0.8181  & 0.8205 $\pm$ 0.0175\\
    \bottomrule
  \end{tabular}
\end{table}

\noindent $^\dag$ Precision/recall values from 5-fold CV are averaged per-fold and may vary by threshold.

\subsection{Summary of Findings}
\begin{itemize}[leftmargin=*]
  \item The paper's AUC of 1.00 is a \textcolor{red}{\textbf{mathematical impossibility}} given the published confusion matrix.
  \item Honest reproduction gives AUC $\approx 0.82$, consistent with published Random Forest literature on Pima.
  \item Accuracy drops from claimed 83.3\% to reproduced 77.9\% under a fresh random seed --- suggesting the paper's 83.3\% may be a best-of-many-runs result or overfit to a favourable split.
\end{itemize}

%---------------------------------------------------------------
\section{Conclusion}
%---------------------------------------------------------------
DiaBeats presents a practical and potentially useful web deployment framework for diabetes risk prediction. The underlying RF model achieves a respectable AUC of $\approx 0.82$ on Pima. However, the paper contains a \textbf{mathematically impossible claim} (AUC = 1.00 with 25 errors), rendering the headline performance numbers untrustworthy. The deployment architecture (Flask/Streamlit + adjustable threshold) is the genuine contribution of this work.

\appendix
\section{Mathematical Proof of AUC Impossibility}
Let $P$ = number of positive examples in the test set = TN + FP + FN + TP, $N$ = negative examples. For the confusion matrix (TN=90, FP=10, FN=15, TP=35):
\begin{align}
  P &= 50,\quad N = 100 \\
  \text{ROC-AUC} &= \frac{1}{PN}\sum_{i: y_i=1}\sum_{j: y_j=0} \mathbf{1}[f(x_i) > f(x_j)]
\end{align}
If AUC = 1, then $f(x_i) > f(x_j)$ for all 50$\times$100 = 5000 pairs. But the classifier misclassifies 10 negatives (FP) at the reported threshold, implying $f(x_j) > 0.5$ for 10 negative examples, and 15 positives (FN) with $f(x_i) \leq 0.5$. For any of these 150 pairs, $f(x_i) \leq f(x_j)$, contradicting AUC = 1. QED.

\section{Artifacts}
\begin{itemize}
  \item Notebook: \texttt{notebooks/03\_kotha\_diabeats.ipynb}
  \item Metrics CSV: \texttt{results/paper03\_kotha\_metrics.csv}
  \item Figure: \texttt{results/fig\_paper03\_diabeats.png}
\end{itemize}

\end{document}
